% ===================================================================
%  ТАБЛИЦЯ № 16 — Великий магазин. — Базар. — Іграшки.
%  Traduction 2026 d'apres texto/io, controlee sur texto/fr, texto/en
%  et texto/es. Consigne : voir texto/uk/10-tablycia-01.tex.
%
%  Le tableau d'astronomie (bloc 15-1) porte des renvois a LETTRES,
%  (a) a (f), graves sur le tableau lui-meme : ils se rangent sous le
%  numero 85, comme le dit gravuri/literi.json. Ce bloc est aussi le
%  troisieme ecart declare de renvoji.py.
%
%  DERNIER TABLEAU DE LA COLONNE UKRAINIENNE. Bilan des seize : zero
%  inversion, comme en catalan, en occitan, en esperanto et en
%  interlingua. La raison est la meme dans les cinq cas, et elle n'a
%  rien a voir avec la parente des langues : toutes postposent le
%  complement du nom, comme l'ido. Deux langues romanes, deux langues
%  construites et une langue slave se rangent ensemble par un seul
%  trait de syntaxe.
% ===================================================================

%%K t16-apar-1 apar
\VUsaut{4.0mm}
\VUtitre{15.0pt}{60}{ТАБЛИЦЯ № 16}
\VUsaut{2.0mm}
\VUpk{13.2pt}{40}{Великий магазин. --- Базар.}
\VUpk{13.2pt}{40}{Іграшки.}
\VUsaut{3.4mm}

%%K t16-01-1 p
1. --- Наближається новий рік. Великі магазини приготували свої виставки
\VUgras{іграшок} \textsuperscript{(1)} і новорічних подарунків.

%%K t16-01-2 p
Я попросив дідуся звести мене на базар подивитися на цю прекрасну виставку, і він
погодився.

%%K t16-02-1 p
2. --- Спершу ми спинилися перед гарними збройними панно, на яких були
\VUgras{рушниця}, \VUgras{кашкет}, \VUgras{шабля}, \VUgras{портупея} і пара
\VUgras{еполетів}, або ж \VUgras{шолом}, \VUgras{кіраса} \textsuperscript{(3)} і
\VUgras{пістолет} \textsuperscript{(4)}. Усе це займало мене куди більше, ніж
\VUgras{чарівні ліхтарі} \textsuperscript{(2)}.

%%K t16-tit-1 sub
\VUsaut{3.0mm}
\VUpk{10.2pt}{40}{Прекрасні видовища.}
\VUsaut{2.6mm}

%%K t16-03-1 p
3. --- У тому самому відділі стояв \VUgras{вартовий} \textsuperscript{(5)} у
своїй \VUgras{будці}; він немов стеріг цілий картонний звіринець. Скільки там
було тварин! \VUgras{Папуга} \textsuperscript{(6)} на жердинці,
\VUgras{носоріг} \textsuperscript{(7)} з рогом на носі, \VUgras{північний олень}
\textsuperscript{(8)}, \VUgras{страус} \textsuperscript{(9)}, \VUgras{мавпа}
\textsuperscript{(10)}, що гризе горіх, \VUgras{лисиця} \textsuperscript{(11)}, що
несе курку, \VUgras{крокодил} \textsuperscript{(12)}, який роззявив величезну
пащу, \VUgras{верблюд} \textsuperscript{(13)}, зграбний \VUgras{хорт}
\textsuperscript{(14)}, \VUgras{пантера} \textsuperscript{(15)}, а потім два
звірі, яких я не знав і які, як мені сказали, --- \VUgras{ласка}
\textsuperscript{(16)} і \VUgras{гієна} \textsuperscript{(17)}. Був там ще
\VUgras{леопард} \textsuperscript{(18)} і \VUgras{вовк} \textsuperscript{(21)};
зовсім поруч --- премилі \VUgras{щури} \textsuperscript{(19)} і крихітні гумові
\VUgras{миші} \textsuperscript{(20)}. Нарешті, \VUgras{бегемот}
\textsuperscript{(22)} і горбатий \VUgras{дромадер} \textsuperscript{(23)}
довершували збірку. Я простояв би там цілі години, якби поруч не опинився чудовий
табір, повний \VUgras{олов'яних солдатиків} \textsuperscript{(29)}, який
привернув мою увагу.

%%K t16-tit-2 sub
\VUsaut{4.0mm}
\VUpk{10.2pt}{40}{Солдати і війна.}
\VUsaut{2.8mm}

%%K t16-04-1 p
4. --- Дідусь, який \VUgras{інвалід} \textsuperscript{(24)} і мусив лишити службу
через \VUgras{рану}, що принесла йому \VUgras{військову медаль}, любить
розповідати мені про \VUgras{походи}, у яких брав участь. Він залюбки пояснив
мені різні пересування цього маленького війська в мініатюрі. Гурт справжніх
солдатів, які оглядали базар --- \VUgras{сапер} \textsuperscript{(25)},
\VUgras{альпійський стрілець} \textsuperscript{(26)} у своєму \VUgras{береті},
піхотний \VUgras{фельдфебель} \textsuperscript{(27)} і артилерійський
\VUgras{сержант} \textsuperscript{(28)} із золотою \VUgras{нашивкою} на рукаві ---
спинився біля нас послухати пояснення.

%%K t16-05-1 p
5. --- Дідусь, дуже гордий таким поважним слухачем, тут же зімпровізував цілий
план битви.

%%K t16-05-2 p
Фортецю з її \VUgras{валами} і \VUgras{ровами} обложило \VUgras{вороже
військо}, яке після довгого \VUgras{бомбардування} готувалося взяти її
\VUgras{приступом}. Але \VUgras{обложені}, доведені голодом, учинили
\VUgras{вилазку} проти \VUgras{табору} \VUgras{обложників}. Відбивши
\VUgras{напад} у жорстокому \VUgras{бою} навколо \VUgras{наметів},
\VUgras{переможні} обложники послали свої \VUgras{ескадрони} в \VUgras{погоню}
за \VUgras{переможеними} нападниками. Ті \VUgras{відступили}, лишивши
\VUgras{поле бою} всіяним \VUgras{убитими} і \VUgras{пораненими}.
\VUgras{Санітари} винесли поранених до \VUgras{польового шпиталю}, щоб їх
лікував \VUgras{хірург}.

%%K t16-05-3 p
Попри геройський опір кількох \VUgras{батальйонів}, що вишикувалися в
\VUgras{каре}, \VUgras{поразка} була цілковита; решта війська \VUgras{втекла},
лишивши безліч \VUgras{полонених}. Ворог готувався увінчати свою
\VUgras{перемогу} взяттям фортеці. Можливо, то був би кінець походу, і після
\VUgras{перемир'я} можна було б підписати \VUgras{мирний договір}.

%%K t16-tit-3 sub
\VUsaut{4.4mm}
\VUpk{10.2pt}{40}{Новорічні подарунки.}
\VUsaut{2.8mm}

%%K t16-06-1 p
6. --- Молоді солдати, які сміялися, слухаючи барвисту розповідь ветерана,
поставили йому ще кілька запитань. Що до мене, я обійшов прилавок, щоб
подивитися на прекрасний \VUgras{арбалет} \textsuperscript{(32)} і на
\VUgras{лук} \textsuperscript{(33)} з його \VUgras{стрілою}
\textsuperscript{(31)}. Там я знайшов ще одну збірку тварин: двох \VUgras{співочих
птахів} \textsuperscript{(34)}, накрутних \VUgras{солов'я} і
\VUgras{кропив'янку}, що співали на все горло, і низку дивних створінь ---
\VUgras{кажана} \textsuperscript{(35)}, \VUgras{удава} \textsuperscript{(36)},
\VUgras{черепаху} \textsuperscript{(37)}, \VUgras{тюленя} \textsuperscript{(38)},
\VUgras{кита} \textsuperscript{(39)} і \VUgras{акулу} \textsuperscript{(40)} з
бодрюшу.

%%K t16-07-1 p
7. --- Я недовго на них дивився, бо помітив те, про що мріяв: чудовий
\VUgras{ляльковий театр} \textsuperscript{(41)} з прекрасними
\VUgras{декораціями} і чарівною червоною \VUgras{завісою}. На \VUgras{сцені} два
актори немов грали \VUgras{роль} у найцікавішій \VUgras{п'єсі}, бо маленькі
картонні \VUgras{глядачі} в ложах, а також ті, що сиділи в \VUgras{партері} і на
\VUgras{гальорці} позаду \VUgras{диригента}, гаряче плескали.

%%K t16-07-2 p
На жаль, той театр був дуже дорогий.

%%K t16-08-1 p
8. --- До того ж вибрати серед іграшок було важко, бо вітрини були завалені
іграшками всякого роду: \VUgras{Арлекін} \textsuperscript{(42)},
\VUgras{Полішинель} \textsuperscript{(43)}, \VUgras{П'єро} \textsuperscript{(44)},
\VUgras{сови} \textsuperscript{(45)}, що плескають крилами, свистючі
\VUgras{дрозди} \textsuperscript{(46)}, гавкучі \VUgras{пуделі},
\VUgras{грамофон} \textsuperscript{(47)} і \VUgras{головоломки}. Одна з цих
іграшок дуже мене потішила: вона зображувала \VUgras{морський бій}
\textsuperscript{(48)}, у якому \VUgras{корабель} атакований \VUgras{піратами}
коло берега, засадженого \VUgras{кокосовими пальмами}.

%%K t16-noto-1 noto
\VUnotes{143.2mm}{%
\textit{(*) Див. таблицю № 6.}%
}

%%K t16-09-1 p
9. --- Усі ці дива я міг розглядати зовсім спокійно, бо в магазині було мало
народу --- ледве жменька роззяв, \VUgras{драгун} \textsuperscript{(49)} та
\VUgras{інкасатор} \textsuperscript{(50)}, який подавав \VUgras{вексель} у
головній \VUgras{касі} \textsuperscript{(51)}.

%%K t16-10-1 p
10. --- Я спитав дідуся, нащо книжки на полицях позаду \VUgras{касира}; він
сказав мені, що це \VUgras{конторські книги}.

%%K t16-tit-4 sub
\VUsaut{3.6mm}
\VUpk{10.2pt}{40}{Оглядаючи магазин.}
\VUsaut{2.8mm}

%%K t16-11-1 p
11. --- Вийшовши з відділу іграшок, ми пішли до лимарського, щоб глянути на
\VUgras{сідла} \textsuperscript{(53)}, \VUgras{стремена} \textsuperscript{(54)},
\VUgras{вуздечки} \textsuperscript{(55)}, \VUgras{вудила} \textsuperscript{(56)}
і \VUgras{хлисти}. Поки \VUgras{завідувач відділу} \textsuperscript{(57)} давав
дідусеві довідки, я пішов дивитися виставку \VUgras{порцеляни} і \VUgras{скла}
\textsuperscript{(58)}, де були прекрасні \VUgras{салатниці} і зграбні
\VUgras{чайники} \textsuperscript{(59)}.

%%K t16-12-1 p
12. --- Щоб піднятися на перший поверх, ми пройшли позаду вітрини
\VUgras{корсетів} \textsuperscript{(60)}, поруч із панчішним відділом, де були
розкладені прекрасні \VUgras{панталони} \textsuperscript{(63)}. Поруч
\VUgras{продавщиця} \textsuperscript{(61)} помагала \VUgras{покупчині}
\textsuperscript{(62)} вибирати прекрасні шовкові \VUgras{тканини}
\textsuperscript{(64)}.

%%K t16-12-2 p
Підіймаючись сходами, я помітив, що магазин оздоблено японськими
\VUgras{ліхтариками} \textsuperscript{(65)}, \VUgras{парасольками}
\textsuperscript{(66)} і величезними \VUgras{пальмами} \textsuperscript{(70)}.

%%K t16-13-1 p
13. --- На першому поверсі дивитися було майже нічого: тільки меблі (*),
\VUgras{вогнетривкі шафи} \textsuperscript{(67)}, \VUgras{комоди}
\textsuperscript{(68)} і \VUgras{дитячі візочки} \textsuperscript{(69)}. Зате
загальний вигляд магазину був дуже \VUgras{цікавий}.

%%K t16-14-1 p
14. --- Під нами я бачив музичні інструменти: \VUgras{арфи}
\textsuperscript{(71)}, \VUgras{гітари} \textsuperscript{(72)},
\VUgras{мандоліни} \textsuperscript{(73)}, \VUgras{корнети}
\textsuperscript{(74)}, \VUgras{кларнети} \textsuperscript{(75)}, скрипки з
їхніми \VUgras{смичками} \textsuperscript{(76)} і навіть \VUgras{великий барабан}
\textsuperscript{(77)}. Трохи далі, коло паперового прилавка, \VUgras{студент}
\textsuperscript{(78)} торгувався з \VUgras{продавцем} \textsuperscript{(79)} про
теку на папери. Біля них я розрізняв прекрасну \VUgras{абетку}
\textsuperscript{(80)}.

%%K t16-14-2 p
Посеред магазину стояла висока \VUgras{різдвяна ялинка} \textsuperscript{(81)},
уся обвішана свічками, дрібничками та іграшками всякого роду:
\VUgras{дерев'яними кониками} \textsuperscript{(82)}, \VUgras{ляльками}
\textsuperscript{(83)} тощо.

%%K t16-14-3 p
\VUgras{Парфумерний відділ} \textsuperscript{(84)} зі своїми \VUgras{зубними
еліксирами} і різними \VUgras{парфумами} видавав дуже приємний запах, який
підіймався до нас.

%%K t16-tit-5 sub
\VUsaut{3.4mm}
\VUpk{10.2pt}{40}{Ми дещо купуємо.}
\VUsaut{2.8mm}

%%K t16-15-1 p
15. --- Що дідусеві ставало нудно, ми знову спустилися великими сходами. Він на
хвилину спинився перед \VUgras{астрономічною таблицею} \textsuperscript{(85)}, на
якій, як він мені сказав, були зображені \VUgras{комети} \textsuperscript{(a)},
\VUgras{затемнення місяця і сонця} \textsuperscript{(b)}, роза вітрів із
\VUgras{чотирма сторонами світу} \textsuperscript{(c)} \VUgras{(північ, південь,
схід і захід)}, карта двох \VUgras{півкуль} \textsuperscript{(d)},
\VUgras{планети} \textsuperscript{(e)} і \VUgras{сонячна система}
\textsuperscript{(f)}. Я нічого в цьому не зрозумів, і мене куди більше займала
галунна ліврея \VUgras{посильного} \textsuperscript{(86)}, який проходив повз, і
гарненькі \VUgras{віяла} \textsuperscript{(87)} поруч зі мною, коло
\VUgras{гаманців} \textsuperscript{(88)} і \VUgras{портмоне}
\textsuperscript{(89)}.

%%K t16-16-1 p
16. --- Я милувався також на прилавку \VUgras{пір'ям} \textsuperscript{(90)} і
\VUgras{пряжками до поясів} \textsuperscript{(91)}, і гарненькими \VUgras{штучними
квітами} \textsuperscript{(94)}, зграбно вкладеними в кошики. \VUgras{Фіалки},
\VUgras{левкої}, \VUgras{гіацинти} \textsuperscript{(95)}, \VUgras{герані}
\textsuperscript{(96)}, \VUgras{лілії} \textsuperscript{(97)} і
\VUgras{настурції} \textsuperscript{(98)} були дуже добре підроблені.

%%K t16-tit-6 sub
\VUsaut{4.4mm}
\VUpk{10.2pt}{40}{Дідусів подарунок.}
\VUsaut{2.8mm}

%%K t16-17-1 p
17. --- Я попросив дідуся купити мені одну із \VUgras{зубних щіток}
\textsuperscript{(92)} у коробці біля \VUgras{шпильок} \textsuperscript{(93)}. Він
погодився, і я сам пішов заплатити за свою покупку в касу, коло якої
зав'язували великий \VUgras{пакунок} \textsuperscript{(100)} для одного
\VUgras{покупця} \textsuperscript{(99)}.

%%K t16-18-1 p
18. --- Раптом серед людей, що стояли коло мистецьких \VUgras{бронз}
\textsuperscript{(101)}, знявся легкий заколот. \VUgras{Злодія}
\textsuperscript{(102)} щойно спіймав на гарячому \VUgras{магазинний детектив}
\textsuperscript{(103)} тієї хвилини, коли той намагався когось обікрасти.
Послали по городового, щоб його \VUgras{затримати}, і якраз коли ми виходили,
винного вели до дільниці.
\VUsaut{6.0mm}
\VUfilet{22mm}
